\section{Where This Project's Model Is Stronger}

\begin{itemize}
  \item \textbf{Full wave-optics blur, not an approximation.} Shen's own
  Discussion names the RSS blur combination as an engineering
  approximation and states that a faithful treatment ``requires the full
  defocused optical transfer function'' --- exactly the upgrade path this
  project already completed via Angular Spectrum propagation.
  \item \textbf{Ghost formula validated against wave optics, not just ray
  tracing.} $R^2=0.995$ against direct two-hole coherent propagation with
  peak detection, versus a formula checked against matrix optics and
  numerical ray tracing (a real but lower-fidelity validation standard for
  a diffractive effect).
  \item \textbf{Gaze-integrated vignetting.} No equivalent exists in
  Shen's model. This project quantified that ignoring gaze angle
  overstates safety for the majority of designs in several families.
  \item \textbf{Margin-aware gates, not boundary-adjacent pass/fail.}
  Following several found-and-fixed defects (blackout designs passing at
  zero margin, a density formula favoring some families by up to 2.25x),
  the final gates require real headroom above a threshold, not merely a
  positive sign.
\end{itemize}

\section{Where Shen's Model Is Stronger, or This Project Has Gaps}

\begin{itemize}
  \item \textbf{Aperture shape.} Shen tests circle, square, diamond,
  ellipse and annulus under a common fill-factor budget and derives a
  rigorous isoperimetric argument for why circular apertures minimize
  blur and ghost rate simultaneously. This project has not tested
  non-circular apertures at all --- circular was adopted by scope
  decision, not by comparison.
  \item \textbf{Closed-form blind-rate.} Shen's $1-\pi/4$ (square) and
  $1-\pi/(2\sqrt3)$ (hex) geometric floors are exact and were not part of
  this project's own metric until very recently, and are still not
  computed for the 8 non-lattice families tested here.
  \item \textbf{Broader, faster search.} The neural surrogate
  ($R^2=0.979$--$0.995$) and inverse-design network let Shen explore many
  more design conditions per unit compute time. This project's search is
  a direct sweep; broader parameter ranges (e.g.\ finer pitch resolution
  near each family's own packing limit) were identified as gaps but not
  fully closed due to time constraints.
  \item \textbf{Perceptual fusion calibration.} Shen's $c_{\min}$
  parameter, though acknowledged as only qualitatively matching clinical
  diplopia onset, gives ghosting a real perceptual floor grounded in
  cited clinical data. This project's hard gate has no equivalent
  perceptual-floor concept --- a margin of exactly zero and a margin of
  one micrometer are treated identically as ``unsafe.''
\end{itemize}

\honestbox{The comparison in Sections 3--4 shows Shen's \emph{ghost} term
specifically fails to discriminate on this project's data. It does not
show Shen's fill-factor, coverage, or aperture-shape analysis is wrong ---
those parts of the published model are untouched by this finding and
remain useful, independently-derived reference points.}

\newpage
